\documentclass[11pt]{article}
\usepackage[margin=1in]{geometry}
\usepackage[T1]{fontenc}
\usepackage[utf8]{inputenc}
\usepackage{lmodern}
\usepackage{microtype}
\usepackage{amsmath,amssymb,amsthm,mathtools,bm}
\usepackage{graphicx}
\usepackage{booktabs}
\usepackage{longtable}
\usepackage{enumitem}
\usepackage{xcolor}
\usepackage[most]{tcolorbox}
\usepackage{hyperref}
\hypersetup{colorlinks=true,linkcolor=blue,citecolor=blue,urlcolor=blue}

\newtcolorbox{insight}[1][]{colback=blue!5!white,colframe=blue!55!black,title={\bfseries Key Insight#1},fonttitle=\bfseries}
\newtcolorbox{formulabox}[1][]{colback=green!5!white,colframe=green!45!black,title={\bfseries #1},fonttitle=\bfseries}
\newtcolorbox{warningbox}[1][]{colback=red!5!white,colframe=red!55!black,title={\bfseries #1},fonttitle=\bfseries}
\newtcolorbox{goalbox}[1][]{colback=green!4!white,colframe=green!40!black,title={\bfseries #1},fonttitle=\bfseries}

\newcommand{\task}[1]{\medskip\noindent\textbf{Task #1:}\ }

\title{Research Plan: Topological Quantum Field Theory,\\
Generalized Symmetries, and Topological Order\\
{\normalsize A Pedagogical Review Bridging High-Energy and Condensed-Matter Viewpoints}}
\author{Helena Brittain and Brian Siew}
\date{\today}

\begin{document}
\maketitle

\begin{abstract}
This document records the plan of attack for the remainder of our Physics 253b final
paper on topological quantum field theory, generalized symmetries, and topological order.
The chapter on Chern--Simons theory (\texttt{paper/chern\_simons\_theory.tex}) has been
drafted and serves as the mathematical anchor of Part~I. The present plan fixes the scope
and author split, enumerates the writing tasks required to complete both halves of the
paper, identifies potential pitfalls, and lays out a nine-day schedule leading into
the final deadline and the one-hour blackboard presentation.
\end{abstract}

\tableofcontents
\newpage

\section{Project Overview}\label{sec:overview}

\begin{goalbox}[Paper Goal]
Write a pedagogical review of topological quantum field theory that (i) is comprehensible
to a graduate student who has completed Schwartz \cite{Schwartz2014QFT} through Chapter~30,
(ii) derives all nontrivial mathematical machinery in full, (iii) grounds the theory in
experimentally accessible physics via the fractional quantum Hall effect, anyon
interferometry experiments, and topological quantum computation, and (iv) meets the
course rubric, in which editorial taste and judgment are the primary grading axes.
\end{goalbox}

\subsection{Core Ingredients Required}

\begin{enumerate}[label=(\roman*)]
\item A mathematical buildup of differential forms, Lie-algebra-valued connections, gauge
transformations, and the Atiyah--Segal axiomatic framework, developed in full for a
reader who has not previously encountered this language (Part~I).
\item Three canonical worked examples, shown end-to-end with explicit calculations:
2D Frobenius TQFT, Chern--Simons theory, and Dijkgraaf--Witten theory (Part~I).
\item A focused treatment of generalized symmetries and anomaly inflow, with the high-energy
chiral-anomaly example (from Schwartz, Chapter~30) and the condensed-matter SPT-edge
example drawn as two instances of the same formal structure (Part~I).
\item Topological order as the infrared description of gapped phases, developed
through the toric code and fractional quantum Hall realizations (Part~II).
\item Anyons, modular tensor category data, and the connection to Chern--Simons
Wilson-loop correlators, sourced directly from Part~I (joint).
\item A careful, honest treatment of the experimental status of anyons: Nakamura et al.\
2020 \cite{Nakamura2020}, Werkmeister et al.\ 2024, Andersen et al.\ 2023, including
the distinction between natural topological phases and quantum-processor emulations
(Part~II).
\item Topological quantum computation, with an explicit statement of which parts of the
topological-protection story are achieved in principle versus in current hardware (Part~II).
\item A unifying outlook chapter tying Parts~I and II together (joint).
\end{enumerate}

\subsection{Project Milestones}

\begin{description}
\item[M1:] Outline v3 locked and repository skeleton set up, with author split confirmed
section by section.
\item[M2:] Part~I (Chapters 2--7) draft complete, including the already-written
Chern--Simons chapter.
\item[M3:] Part~II (Chapters 8--12) draft complete.
\item[M4:] Front matter (abstract, introduction), back matter (outlook, bibliography),
and all appendices drafted.
\item[M5:] Full paper internally cross-referenced, figures drafted or sourced, bibliography
consolidated, typeset clean, and proofread end-to-end.
\item[M6:] Blackboard-and-chalk presentation rehearsed to time with the intended board
layout and handoff between authors.
\end{description}

\subsection{Standing Assumptions}

\begin{insight}[Editorial Principle]
The course rubric grades taste and judgment more heavily than completeness. We will
\emph{deliberately} cut candidate sections (non-invertible symmetries, SymTFT in depth,
Reshetikhin--Turaev in detail, modular tensor category pentagon/hexagon, full Jones-polynomial
derivation via surgery) to keep the paper's through-line clear.
\end{insight}

Our through-line, which will appear in the abstract and again in the outlook, is:
\emph{Topological quantum field theory is the infrared language of gapped phases. We
demonstrate this on three touchstones---Frobenius 2D TQFT, Chern--Simons theory, and
Dijkgraaf--Witten theory---and use them to motivate generalized symmetries, anomaly inflow,
and the experimental anyon program.} Every chapter must tie back to one of these
touchstones or be cut.

\subsection{Author Split}\label{sec:split}

\begin{itemize}
\item \textbf{Brian (B):} mathematical and axiomatic side. Differential forms primer,
axiomatic TQFT, 2D Frobenius TQFT, Chern--Simons (already drafted), BF theory,
Dijkgraaf--Witten, generalized symmetries, anomaly inflow, modular tensor category
formalism, and the mathematical appendices.
\item \textbf{Helena (H):} physical and experimental side. Topological order as
infrared TQFT, fractional quantum Hall effect, anyon-interferometry experiments,
topological quantum computation, and the experimental honesty boxes.
\item \textbf{Joint (B+H):} introduction, abstract, outlook, modular data chapter
(Brian on formalism, Helena on physics), figure captions, final polish.
\end{itemize}

\section{Paper Structure and Scope}\label{sec:structure}

\subsection{Part I: Theoretical Buildup}

\begin{enumerate}[leftmargin=2em]
\item Introduction and scope (joint, $\sim 6$~pp)
\item Differential geometry primer: forms, gauge theory, bundles (B, $\sim 10$~pp)
\item Axiomatic TQFT: cobordisms and the Atiyah--Segal axioms (B, $\sim 10$~pp)
\item 2D TQFT and commutative Frobenius algebras --- mathematical showpiece (B, $\sim 12$~pp)
\item \textbf{Chern--Simons theory} (B, \emph{already drafted})
\item BF theory and the bridge to Dijkgraaf--Witten (B, $\sim 8$~pp)
\item Dijkgraaf--Witten theory: finite-group gauge theories and group cohomology (B, $\sim 10$~pp)
\item Generalized symmetries and anomaly inflow (B, $\sim 14$~pp)
\end{enumerate}

\subsection{Part II: Physical Applications and Experiment}

\begin{enumerate}[leftmargin=2em,start=9]
\item Topological order as infrared TQFT (H, $\sim 8$~pp)
\item Anyons, braiding, fusion, and modular data (B+H, $\sim 10$~pp)
\item Fractional quantum Hall effect as a physical realization (H, $\sim 10$~pp)
\item Experiments on anyonic braiding (H, $\sim 10$~pp)
\item Topological quantum computation (H, $\sim 8$~pp)
\item Outlook and synthesis (joint, $\sim 4$~pp)
\end{enumerate}

\subsection{Appendices}

\begin{enumerate}[label=\Alph*.,leftmargin=2em]
\item Category theory in 6 pages (B)
\item Frobenius algebras and the 2D classification in detail (B)
\item Chern--Simons level quantization (B, sourced from the CS chapter)
\item BF theory and $\mathbb{Z}_N$ gauge theory worked out (B)
\item Topological entanglement entropy derivation (B)
\end{enumerate}

\begin{insight}[Decisions Already Made]
(i)~No non-invertible symmetry section; one paragraph in the outlook instead.
(ii)~SymTFT compressed to a single page inside the generalized-symmetries chapter.
(iii)~No full derivation of the Jones polynomial; the Chern--Simons chapter states the
result and cites Witten.
(iv)~No Reshetikhin--Turaev categorical details; the TQFT axioms chapter states the theorem.
\end{insight}

\section{Stage 0: Precondition}\label{sec:stage0}

\noindent\textit{Status: complete. Summary for reference.}

\begin{itemize}
\item \texttt{paper/chern\_simons\_theory.tex} is drafted at $\sim 25$~pp with full
derivations of the algebraic identity $\tr(F\wedge F)=d\omega_{\CS}$, the gauge transformation
law for $\omega_{\CS}$, the Euler-angle evaluation of $\int_{SU(2)}\tr((g^{-1}dg)^{\wedge 3})=24\pi^2$,
and the finite Heisenberg-algebra derivation of $\dim\mathcal{H}_{U(1),k}(T^2)=k$.
\item \texttt{paper/chern\_simons\_theory.bib} contains all citations used by the CS chapter.
\item Outline v2 (\texttt{TQFT\_outline\_v2.tex}) exists; v3 will be produced in Stage~I.
\item Style references \texttt{hdr.tex}, \texttt{springer.tex} have been read and their
conventions (\texttt{\textbackslash defthm}, trace normalization, preferred environments)
are to be matched in all new chapters.
\end{itemize}

\section{Stage I: Scope Freeze and Repository Skeleton}\label{sec:stageI}

\begin{goalbox}[Objective]
Lock the outline at v3, create the full repository skeleton with placeholder \texttt{.tex}
files for every chapter, consolidate a single master bibliography, and confirm the author
split at section granularity. Deliverables must be shareable with Helena in one click.
\end{goalbox}

\subsection{Outline v3 and Author Split}

\task{1.1} Produce \texttt{TQFT\_outline\_v3.tex} by merging outline v2 with the
chapter list of Section~\ref{sec:structure} and the author split of
Section~\ref{sec:split}. Mark each chapter \textsc{Deep}, \textsc{Bridge}, or \textsc{Survey};
assign initial page budget.

\task{1.2} Walk Helena through outline v3 and confirm her ownership of Chapters 9, 11,
12, and 13 as well as the experimental honesty boxes. Record any disagreements in
\texttt{plan/notes/author\_split\_decisions.md}.

\task{1.3} Decide the placement of the modular tensor category formalism. Default: the
formalism subsection in Chapter~10 is Brian's; the physics framing is Helena's. Alternatives
recorded if disagreed.

\subsection{Repository Skeleton}

\task{1.4} Create \texttt{paper/main.tex} with \texttt{\textbackslash include}-style
structure calling \texttt{paper/ch01\_intro.tex}, \texttt{paper/ch02\_forms.tex},
\texttt{paper/ch03\_axiomatic.tex}, \texttt{paper/ch04\_frobenius.tex},
\texttt{paper/chern\_simons\_theory.tex} (already present), \texttt{paper/ch06\_bf.tex},
\texttt{paper/ch07\_dw.tex}, \texttt{paper/ch08\_gensym.tex},
\texttt{paper/ch09\_toporder.tex}, \texttt{paper/ch10\_anyons.tex},
\texttt{paper/ch11\_fqh.tex}, \texttt{paper/ch12\_experiments.tex},
\texttt{paper/ch13\_tqc.tex}, \texttt{paper/ch14\_outlook.tex}, and appendix files.

\task{1.5} Create a single \texttt{paper/tqft\_review.bib} by merging the CS chapter
bib with outline v2's bibliography and the additional citations we have identified
(Nakamura 2020, Werkmeister 2024, Andersen 2023, Kitaev 2003, Kitaev--Preskill 2006,
Levin--Wen 2006, Freed--Moore--Teleman 2024, Bhardwaj et al.~2023 lectures, GKSW 2015).

\task{1.6} Write a shared preamble file \texttt{paper/preamble.tex} that defines the
macros and theorem environments used in \texttt{hdr.tex} and \texttt{springer.tex},
removing the \texttt{upshot}, \texttt{note}, \texttt{slogan} etc. environments that
we have already excluded from the CS chapter.

\task{1.7} Commit a \texttt{paper/figures/} directory with a \texttt{README.md}
listing the target figure set (see Section~\ref{sec:figures}).

\subsection{Figure and Table Plan}\label{sec:figures}

\task{1.8} Commit a checklist of every figure in the paper, marking \textsc{generate},
\textsc{extract (with citation)}, or \textsc{no figure}. Initial list:
\begin{itemize}[leftmargin=2em]
\item Cobordism generators (cap, cup, cylinder, pair of pants) --- generate
\item Pair-of-pants multiplication/comultiplication --- generate
\item Wilson-loop link diagram (Hopf link) --- generate
\item Higher-form symmetry defect insertion --- generate
\item Bulk--boundary anomaly inflow schematic --- generate
\item Anyon braiding worldlines --- generate
\item FQH interferometer schematic --- extract from Nakamura 2020 with citation
\item Fusion tree / Ising braiding diagram --- generate
\end{itemize}

\task{1.9} Commit a checklist of every table: the generalized-symmetry dictionary table,
the anyon-comparison table (toric code / Laughlin / Ising), and the page-budget table.

\section{Stage II: Part I Mathematical Buildup A}\label{sec:stageII}

\begin{goalbox}[Objective]
Write Chapters 2, 3, and 4: the differential-geometry primer, the axiomatic
Atiyah--Segal framework, and the 2D TQFT / Frobenius-algebra showpiece.
\end{goalbox}

\subsection{Chapter 2: Differential Geometry Primer}

\begin{insight}[Pedagogical Stance]
The reader has Schwartz-level facility with components, index notation, and the chiral
anomaly of Chapter~30, but typically no formal exposure to differential forms as
geometric objects, principal bundles, or connections beyond the local Yang--Mills
expression. Chapter~2 bridges this gap without becoming a standalone geometry text.
\end{insight}

\task{2.1} Draft Section~2.1 (forms, wedge, exterior derivative). Reuse
\texttt{paper/chern\_simons\_theory.tex} Section~\ref{sec:forms}: pull the Lemma
\texttt{lem:leibniz-nilpotent} and its proof into Chapter~2 and cite back from
Chapter~5 rather than duplicating.

\task{2.2} Draft Section~2.2 (integration, Stokes, de Rham cohomology). Verify the
Stokes statement against Nakahara and cite without hallucinating a theorem number.

\task{2.3} Draft Section~2.3 (Lie-algebra-valued forms, field strength, anti-Hermitian
convention and its translation to Schwartz). State once, cite forward to Chapter~5 for
all applications.

\task{2.4} Draft Section~2.4 (principal bundles, connections one-form level): minimal
and motivated by the need for $\omega_{\CS}$ to be defined on the total space.

\task{2.5} Draft Section~2.5 (gauge transformations and the transformation law of $F$).
Reuse the template computation from the CS chapter verifying $F^g=g^{-1}Fg$.

\task{2.6} Include two worked examples: electromagnetism in forms language, and the
winding number of $S^1\to S^1$ as an integer-valued integral.

\task{2.7} Verify chapter is around 10~pp and feeds Chapter~5 directly; remove anything
not needed by Chapters 3--8.

\subsection{Chapter 3: Axiomatic TQFT}

\task{3.1} Draft Section~3.1 introducing the cobordism category: objects, morphisms,
composition, symmetric monoidal structure. Use one concrete worked example.

\task{3.2} Draft Section~3.2 stating the Atiyah--Segal axioms: $Z(\Sigma_1\sqcup\Sigma_2)=
Z(\Sigma_1)\otimes Z(\Sigma_2)$, gluing $\to$ composition, $Z(\emptyset)=\C$.
Prove no theorem here; state the axioms and derive immediate consequences (finite
dimensionality of $Z(\Sigma)$ from cylinder-as-identity plus trace-finite arguments).

\task{3.3} Draft Section~3.3 explicitly gluing two disks to form a sphere and relating
partition-function normalizations to the cylinder trace. This is the one worked gluing
example demanded by the outline.

\task{3.4} Draft Section~3.4 (taxonomy): Schwarz-type versus Witten-type versus
invertible, one paragraph each.

\task{3.5} Cross-check: every claim in Chapter~3 is either a definition, a worked
example, or cited to Atiyah 1988 \cite{Atiyah1988TQFT}.

\subsection{Chapter 4: 2D TQFT and Frobenius Algebras}

\begin{insight}[Showpiece]
Chapter~4 is Brian's primary mathematical showpiece. Present the classification theorem
with the full proof that the generators and relations of $\mathrm{Cob}_2$ produce the
axioms of a commutative Frobenius algebra, and conversely.
\end{insight}

\task{4.1} Draft Section~4.1 stating the generators of $\mathrm{Cob}_2$ (cap, cup,
cylinder, pair of pants) and the relations among them. Include the pair-of-pants diagram
as Figure~1.

\task{4.2} Draft Section~4.2 proving that a 2D oriented TQFT is determined by a finite
dimensional vector space $V=Z(S^1)$ with a commutative, associative, unital
multiplication $m:V\otimes V\to V$ (from the pair of pants), a trace $\epsilon:V\to\C$
(from the disk), and a nondegenerate pairing (from the cylinder), satisfying the
Frobenius condition. Follow Kock \emph{Frobenius Algebras and 2D TQFTs}; do not
introduce category-theory heavy machinery.

\task{4.3} Worked examples: group algebra of a finite group $G$ with the trace
$\epsilon(g)=\delta_{g,e}$ gives the 2D Dijkgraaf--Witten TQFT on genus-$g$
surfaces, and this pre-figures Chapter~7.

\task{4.4} Section~4.3: the 2D partition function on a genus-$g$ surface is
$Z(\Sigma_g) = \sum_i \lambda_i^{2-2g}$ for $\lambda_i$ eigenvalues of multiplication
by $e_i$ idempotent basis. Derive this via the pair-of-pants/cap gluing, in full.

\task{4.5} Verify that Appendix~B will carry the full proof with all details; Chapter~4
presents the argument at review level with key steps.

\section{Stage III: Part I Mathematical Buildup B}\label{sec:stageIII}

\begin{goalbox}[Objective]
Write Chapters 6, 7, and 8: BF theory, Dijkgraaf--Witten, and generalized symmetries
with anomaly inflow.
\end{goalbox}

\subsection{Chapter 6: BF Theory}

\task{6.1} Section~6.1: action $S = \frac{k}{2\pi}\int B\wedge F$, equations of motion
$F=0$ and $dB=0$, Wilson-line and 't~Hooft-surface operators. Derive the equations of
motion explicitly; do not state them.

\task{6.2} Section~6.2: compact $U(1)$ BF theory at level $N$ on a 3-manifold is
equivalent to $\mathbb{Z}_N$ gauge theory. Derive this by integrating out $B$ as
Lagrange multiplier and working out the residual gauge structure. This is a bridge
computation tying to toric code in Chapter~9.

\task{6.3} Section~6.3: Wilson-line/'t Hooft-surface linking number on closed
3-manifolds. Show that $\langle W(C) T(\Sigma)\rangle \propto e^{2\pi i \mathrm{Lk}(C,\Sigma)/N}$,
paralleling the CS abelian calculation already done in \texttt{paper/chern\_simons\_theory.tex}.

\task{6.4} Section~6.4: ground-state degeneracy of BF on $\Sigma\times\R$ with
$\Sigma$ genus $g$ is $N^{2g}$. Derive via flat-connection moduli space counting, same
methodology as the $U(1)_k$ torus Hilbert space in the CS chapter.

\begin{warningbox}[Convention Check]
Confirm the relative-sign convention between $B\wedge F$ and $F\wedge B$ used in Tong's
notes versus our CS conventions. Choose one and cross-reference the CS chapter.
\end{warningbox}

\subsection{Chapter 7: Dijkgraaf--Witten Theory}

\task{7.1} Section~7.1: finite-group gauge theory, path integral as
$Z(M) = \frac{1}{|G|}\sum_{\phi:\pi_1(M)\to G}\,e^{2\pi i\,\langle\alpha,\phi_*[M]\rangle}$,
where $\alpha\in H^d(BG,U(1))$ is the twisting.

\task{7.2} Section~7.2: classify 3D Dijkgraaf--Witten actions by $H^3(G,U(1))$ and
compute this group explicitly for $G=\mathbb{Z}_N$: obtain $H^3(\mathbb{Z}_N,U(1))=\mathbb{Z}_N$
via the bar resolution, in full.

\task{7.3} Section~7.3: worked example --- the untwisted $\mathbb{Z}_2$ theory on $M$ is
the toric code at long wavelengths. Connect explicitly to Chapter~6 compact BF theory.

\task{7.4} Section~7.4: twisted Dijkgraaf--Witten theory for $G=\mathbb{Z}_N$ as a
laboratory for anomaly inflow: a 4D cocycle on the boundary of a 5D bulk produces
a 't~Hooft anomaly that cannot be gauged without inflow.

\task{7.5} Cross-reference: Chapter~7 will be cited in the torus Hilbert-space
discussion of the anyons chapter, and the anomaly-inflow chapter, and the outlook.

\subsection{Chapter 8: Generalized Symmetries and Anomaly Inflow}

\begin{insight}[Scope Control]
We deliberately trim to 14~pp total. Cover: $q$-form symmetries, gauging, one anomaly
inflow example of each type (HEP and CM), and SymTFT in a single page. Non-invertible
symmetries get one outlook paragraph.
\end{insight}

\task{8.1} Section~8.1: 0-form symmetries reviewed in the language of topological
defects, not charge operators. Transition from Noether to ``symmetry operator on a
codimension-1 manifold''.

\task{8.2} Section~8.2: $q$-form symmetries. Define: charged object (dimension~$q$),
symmetry operator (codimension $q+1$), background field ($(q+1)$-form). Include the
\emph{dictionary table} mandated by the outline, translating among object, dimension,
charged operator, symmetry defect, conserved current, background field.

\task{8.3} Section~8.3: gauging higher-form symmetries, with the worked $\mathbb{Z}_N$
example reused from Chapter~7. Show explicitly that gauging a 1-form symmetry in
electrodynamics produces magnetic monopoles as charged 0-form symmetry operators.

\task{8.4} Section~8.4: anomalies as obstructions to gauging, stated as follows.
A theory has a 't~Hooft anomaly for symmetry $G$ if the partition function in a
background $G$-field is not gauge-invariant but shifts by a local functional of
the background.

\task{8.5} Section~8.5: anomaly inflow.
Argument: a $(d+1)$-dimensional invertible TQFT on a manifold with boundary produces a
non-gauge-invariant shift at the boundary that cancels the boundary anomaly.

\task{8.6} Section~8.6: high-energy example. Take Schwartz's Chapter~30 chiral anomaly
$\partial_\mu j^\mu_A = (1/16\pi^2)\epsilon^{\mu\nu\rho\sigma}\tr F_{\mu\nu}F_{\rho\sigma}$
and reframe it as the boundary of a 5D Chern--Simons bulk. Work through the bulk action
explicitly.

\task{8.7} Section~8.7: condensed-matter example. The edge chiral fermion of a 3D
topological insulator has a mixed anomaly in $U(1)_{\text{charge}}\times T$ that is
canceled by the bulk $\theta$-term at $\theta=\pi$. Do this at a level of detail
comparable to Section~8.6.

\task{8.8} Section~8.8: SymTFT in one page. Slogan, one picture of the three-layer
sandwich, and the DW / toric-code example. No formalism.

\task{8.9} Scope guardrail: at the end of this chapter, confirm the chapter is
$\le 14$~pp. If over, trim Section~8.8 first, then Section~8.4.

\section{Stage IV: Part II Physical Applications}\label{sec:stageIV}

\begin{goalbox}[Objective]
Write Chapters 9, 10, 11: topological order as infrared TQFT, anyons and modular data,
fractional quantum Hall effect. Helena leads Chapters~9 and 11; Chapter~10 is jointly
authored.
\end{goalbox}

\subsection{Chapter 9: Topological Order as Infrared TQFT}

\task{9.1} Section~9.1 (H): long-range entanglement as the modern definition of
topological order. Explicit contrast with symmetry-breaking order.

\task{9.2} Section~9.2 (H): ground-state degeneracy on nontrivial manifolds as a
signature of topological order. Worked example: $\mathbb{Z}_2$ toric code on the torus
gives 4-fold degeneracy, derived from the explicit Hamiltonian
$H = -\sum_v A_v - \sum_p B_p$ and the $W_1, W_2$ string operator algebra.

\task{9.3} Section~9.3 (H): emergent anyons from the string operators of the toric
code. Braiding statistics calculated explicitly.

\task{9.4} Section~9.4 (H, pointing to B): the infrared effective theory of the toric
code is $\mathbb{Z}_2$ BF theory; cite Chapter~6. This is the round-trip back to Part~I.

\subsection{Chapter 10: Anyons, Braiding, Fusion, and Modular Data}

\task{10.1} Section~10.1 (H): abelian vs.\ non-abelian anyons, fusion rules, the
pentagon axiom. Do not write the $F$-symbols.

\task{10.2} Section~10.2 (H): braiding, the hexagon axiom, topological spin $\theta_a$.
No derivation; state the axioms, draw the diagrams.

\task{10.3} Section~10.3 (B): modular $S$ and $T$ matrices. Derive the $S$ matrix for
$U(1)_k$ explicitly from the $k$-dimensional Hilbert space of
\texttt{paper/chern\_simons\_theory.tex} (Proposition~\ref{prop:dimU1k} there).
Show that $S$ diagonalizes the fusion rules.

\task{10.4} Section~10.4 (B): connection to Chern--Simons. For $U(1)_k$ and for
$SU(2)_k$, list the primary labels, $S$ and $T$, and the relation to CS Wilson-loop
correlators already derived in the CS chapter.

\task{10.5} Section~10.5 (H+B): suggested table comparing toric code anyons, Laughlin
anyons, and Ising anyons: labels, fusion, braiding phases, topological spin.

\subsection{Chapter 11: Fractional Quantum Hall Effect}

\task{11.1} Section~11.1 (H): physical setup of the quantum Hall effect. Landau levels,
filling fraction, integer vs.\ fractional.

\task{11.2} Section~11.2 (H): Laughlin wavefunction for $\nu = 1/m$, and the emergence
of $U(1)_m$ Chern--Simons as the effective long-wavelength theory. Derive the
$\sigma_{xy} = e^2/(mh)$ relation from the CS Lagrangian by integrating out $A$ against
a background source and cite back to the CS chapter.

\task{11.3} Section~11.3 (H): quasiparticles as Wilson-line endpoints, $e/m$ charge,
$2\pi/m$ braiding phase. Cite Example~\ref{ex:abelian-wilson} of the CS chapter for
the derivation.

\task{11.4} Section~11.4 (H): edge theory as chiral boson. Anomaly inflow between the
bulk CS term and the edge chiral boson; cite Chapter~8.

\task{11.5} Section~11.5 (H): honesty box distinguishing experimentally established
abelian anyonic braiding from the more tentative non-abelian platforms. This is
required by the outline.

\section{Stage V: Part II Experiments and Applications}\label{sec:stageV}

\begin{goalbox}[Objective]
Write Chapters 12 (experiments) and 13 (topological quantum computation), Helena lead.
\end{goalbox}

\subsection{Chapter 12: Experiments on Anyonic Braiding}

\task{12.1} Section~12.1 (H): Nakamura--Liang--Gardner--Manfra \cite{Nakamura2020}
interferometer at $\nu=1/3$. Principle of the measurement, the $e^{2\pi i/3}$ phase,
identification with the CS Wilson-loop prediction of Chapter~5.

\task{12.2} Section~12.2 (H): Werkmeister et al.\ 2024 graphene interferometer; discuss
improvements and the specific claims.

\task{12.3} Section~12.3 (H): Andersen et al.\ 2023, non-abelian braiding on a
superconducting processor. Important \emph{honesty box}: this is a synthetic
emulation of braiding on hardware, not a natural topological phase. Say so directly.

\task{12.4} Figure: reproduce one figure from each experiment with full credit and
caption in the format of the style references.

\subsection{Chapter 13: Topological Quantum Computation}

\task{13.1} Section~13.1 (H): non-abelian anyons and the promise of fault-tolerance.
Cite Kitaev 2003 for the anyonic quantum-computation proposal.

\task{13.2} Section~13.2 (H): braids as gates and the role of topological degeneracy.
One example (Ising anyons) worked through carefully.

\task{13.3} Section~13.3 (H): honesty box on current experimental limitations: no
natural non-abelian platform has been confirmed, Majorana-based proposals remain unverified,
and synthetic-braiding demonstrations are not yet fault-tolerant in the strict sense.

\section{Stage VI: Integration and Polish}\label{sec:stageVI}

\begin{goalbox}[Objective]
Draft front matter and outlook, consolidate the bibliography, install all figures,
resolve all cross-references, and proofread the full paper.
\end{goalbox}

\subsection{Front Matter}

\task{14.1} Draft the abstract. Keep it under 250 words; mention the through-line
(three touchstones), the ``from Schwartz Chapter~30'' calibration, and the experimental
anchor in anyon interferometry.

\task{14.2} Draft Chapter~1, the introduction. Scope subsection, thesis, target audience,
how to read the paper. Do not summarize every chapter; just say which part covers what.

\task{14.3} Draft Chapter~14, the outlook. Non-invertible symmetries in one paragraph.
SPT/SET classification beyond group cohomology in one paragraph. Fermionic topological
order in one paragraph. Fault-tolerant architectures in one paragraph.

\subsection{Integration}

\task{14.4} Propagate all cross-references: \texttt{\textbackslash ref}s across chapters
must resolve. Run \texttt{latexmk} three times and confirm no ``undefined reference''
warnings.

\task{14.5} Consolidate \texttt{tqft\_review.bib}; ensure every \texttt{\textbackslash cite}
resolves and no bib entry is orphaned.

\task{14.6} Verify every chapter ends with a 1--2 sentence pointer into the next chapter;
this is the pedagogical glue.

\task{14.7} Generate all figures not yet sourced (cobordism diagrams, anyon braiding
diagrams, anomaly inflow schematic). Use TikZ where possible for reproducibility.

\task{14.8} Double-check every one of the four ``must include'' items from outline v2:
(a)~gluing example in Chapter~3, (b)~FQH physical connection in Chapter~5 (already
satisfied in the CS chapter), (c)~generalized-symmetry dictionary table in Chapter~8,
(d)~honesty box in Chapter~12.

\subsection{Proofreading}

\task{14.9} Full read-through pass~1 (Brian): math correctness. Sanity-check every
numerical factor, every sign, every cohomology-group identity. Flag ghost
references (``Chapter~X of Nakahara''--style) that do not exist.

\task{14.10} Full read-through pass~2 (Helena): physical reasonableness, experimental
numbers, and citation attribution.

\task{14.11} Full read-through pass~3 (joint): prose, transitions, sentence-level
editing. Remove AI-tell phrases (``we shall see'', ``in this chapter we have proved'',
``the key statement'').

\section{Stage VII: Presentation and Final Submission}\label{sec:stageVII}

\begin{goalbox}[Objective]
Prepare and rehearse the one-hour blackboard presentation with the structure we agreed
on (25 min Brian mathematical CS, 25 min Helena CS-as-effective-theory and experiments,
10 min Q\&A), and submit the paper.
\end{goalbox}

\subsection{Blackboard Presentation}

\task{15.1} Brian drafts 25-min board script (chalk blocks, equation persistence,
showpiece at minute~10: level quantization from the large-gauge argument).

\task{15.2} Helena drafts 25-min board script (Hall conductance derivation, anyon
braiding, three experiment summaries).

\task{15.3} Both timers at home, individually, at least twice. Blackboard talks run
long on the first pass; cut aggressively.

\task{15.4} One joint dry-run at least 48 hours before the scheduled presentation.
Iterate on the handoff.

\task{15.5} Print a 1-page \emph{cheat card} per author: equation layout on the board,
minute checkpoints, the 3 hardest questions we expect with prepared answers.

\subsection{Final Submission}

\task{15.6} Compile final PDF. Confirm page count between 75 and 100 pages main text
plus 15--25 appendix. If outside, cut.

\task{15.7} Sanity-check submission format against course policy; submit.

\section{Timeline and Deliverables}\label{sec:timeline}

\begin{center}\small
\begin{tabular}{@{}lllp{6cm}@{}}
\toprule
\textbf{Stage} & \textbf{Duration} & \textbf{Day} & \textbf{Key Deliverable} \\
\midrule
0: Precondition & done & --- & CS chapter, outline v2, bib scaffold \\
I: Scope freeze & 0.5 day & 1 & Outline v3, repo skeleton, figure/table plan \\
II: Part I A & 2 days & 1--3 & Chapters 2, 3, 4 drafts \\
III: Part I B & 2 days & 3--5 & Chapters 6, 7, 8 drafts \\
IV: Part II A & 2 days & 5--7 & Chapters 9, 10, 11 drafts \\
V: Part II B & 1.5 days & 7--8 & Chapters 12, 13 drafts \\
VI: Integration & 1 day & 8--9 & Abstract, intro, outlook, figures, proofread \\
VII: Presentation & 1 day & 9 & Board scripts, dry-runs, final submit \\
\bottomrule
\end{tabular}
\end{center}
\smallskip

\noindent Parallelism: during Stages II--III (Brian's heavy writing load), Helena should
be working on Chapters 9, 11, 12 drafts in parallel, so that her Stage IV work is
primarily revision rather than initial drafting. If Brian's Stage III runs long, drop
Section~8.8 (SymTFT) entirely and recover a half-day.

\section{Key Risks and Open Questions}\label{sec:risks}

\begin{enumerate}[leftmargin=2em]
\item \textbf{Chapter 8 scope blow-up.} The generalized-symmetries material could easily
consume a week if we are not disciplined. Mitigation: strict 14-page cap, SymTFT to one
page, non-invertible to outlook paragraph.
\item \textbf{Ghost references.} Our first CS draft cited ``Lemma 1.45 of Freed'' and
``Appendix B of Nakahara'' that do not exist. Mitigation: Task~14.9 is the scheduled
ghost-reference audit; ban bracketed citations like \texttt{\textbackslash cite[Eq.~X]\{Y\}}
unless we have the actual source open.
\item \textbf{Handoff between Chapters 5 and 11.} The CS chapter already states the
abelian Wilson-loop linking-number result; Chapter~11 must cite rather than re-derive.
Mitigation: explicit \texttt{\textbackslash ref} calls; in Task~14.4 verify.
\item \textbf{MTC formalism ownership.} Chapter~10 splits Brian's math and Helena's
physics. Mitigation: Task~1.3 forces an explicit decision.
\item \textbf{Figure generation time.} TikZ for cobordism diagrams is nontrivial.
Mitigation: budget Task~14.7 early; if blocked, hand-draw and scan cleanly.
\item \textbf{Presentation running long.} Blackboard talks always run $\sim 30\%$ over
the first time. Mitigation: Task~15.3 (individual timed runs) before Task~15.4 (joint run).
\item \textbf{The honesty boxes.} Chapter~12 and Chapter~13 must state openly that
topological quantum computation is not yet realized in hardware. Not negotiable; failure
to state this would be a scholarship error.
\end{enumerate}

\section{Complete Task List}\label{sec:tasks}

For reference, the full task list:

\paragraph{Stage I.} 1.1 outline v3; 1.2 author split walkthrough; 1.3 MTC ownership;
1.4 main.tex skeleton; 1.5 merged bibliography; 1.6 shared preamble; 1.7 figures
directory; 1.8 figure checklist; 1.9 table checklist.

\paragraph{Stage II.} 2.1--2.7 Chapter~2 (forms primer);
3.1--3.5 Chapter~3 (axiomatic TQFT);
4.1--4.5 Chapter~4 (2D TQFT and Frobenius algebras).

\paragraph{Stage III.} 6.1--6.4 Chapter~6 (BF theory);
7.1--7.5 Chapter~7 (Dijkgraaf--Witten);
8.1--8.9 Chapter~8 (generalized symmetries and inflow).

\paragraph{Stage IV.} 9.1--9.4 Chapter~9 (topological order);
10.1--10.5 Chapter~10 (anyons and modular data);
11.1--11.5 Chapter~11 (fractional quantum Hall).

\paragraph{Stage V.} 12.1--12.4 Chapter~12 (experiments);
13.1--13.3 Chapter~13 (topological quantum computation).

\paragraph{Stage VI.} 14.1--14.3 front and back matter;
14.4--14.8 integration; 14.9--14.11 proofreading passes.

\paragraph{Stage VII.} 15.1--15.5 presentation; 15.6--15.7 final submission.

\begin{thebibliography}{9}
\bibitem{Schwartz2014QFT}
M.~D.~Schwartz, \emph{Quantum Field Theory and the Standard Model}, Cambridge University
Press, 2014.
\bibitem{Atiyah1988TQFT}
M.~F.~Atiyah, ``Topological quantum field theories,'' \emph{Publ.\ Math.\ IH\'ES}
\textbf{68} (1988), 175.
\bibitem{Nakamura2020}
J.~Nakamura, S.~Liang, G.~C.~Gardner, M.~J.~Manfra, ``Direct observation of anyonic
braiding statistics at the $\nu=1/3$ fractional quantum Hall state,'' \emph{Nature Physics}
\textbf{16} (2020), 931.
\end{thebibliography}

\end{document}
